\documentclass[12pt, a4paper]{article}

% Required Packages
\usepackage{amsmath, amssymb, graphicx}
\usepackage{kotex} % For Korean text
\usepackage{geometry}
\usepackage{enumitem} % [label=...] 옵션을 사용하기 위해 필수적인 패키지입니다.

\graphicspath{{graph/}}
\geometry{a4paper, margin=1in}
\title{CLINIC 13-14}

\begin{document}
\maketitle

\begin{enumerate}
\item 다음 부정적분을 구하여라.
   % itemize 대신 enumerate를 사용하고, enumitem 패키지 덕분에 아래 옵션이 정상 작동합니다.
   \begin{enumerate}[label=(\arabic*)] 
    \item $\displaystyle \int \left( \frac{1}{1+\tan^{2}\theta} + \frac{1}{1+\cot^{2}\theta} \right) d\theta$
    \item $\displaystyle \int (\sin\theta+\cos\theta)^{2}d\theta + \int (\sin\theta-\cos\theta)^{2}d\theta$
    \item $\displaystyle \int \sin^2 \frac{x}{2} dx$
    \item $\displaystyle \int \frac{8^x - 2^x}{2^x + 1} dx$
    \item $\int x(1-x)^{10}dx$
    \item $\int\dfrac{1}{\sqrt[4]{2x+3}}dx$
    \item $\int10^{3x+2}dx$
    \item $\int(e^{x}-e^{-x})^{2}dx$
    \item $\int\dfrac{1}{\sqrt{x+1}-\sqrt{x-1}}dx$
    \item $\int\dfrac{e^{x}}{\sqrt{e^{x}+1}}dx$
    \item $\int\tan 2x~dx$
    \item $\int\dfrac{x^{3}}{\sqrt{1-x^{2}}}dx$
    \item $\int\dfrac{x-1}{\sqrt{x+1}}dx$
    \item $\int x\cos^{2}x~dx$
    \item $\int \ln(x+2)dx$
    \item $\int(\ln x)^{2}dx$
    \item $\int x^{2}e^{x}dx$
    \item $\displaystyle \int\frac{3 \sin^{3}x-3 \sin x+\cos^{3}x-2}{\cos^{2}x}dx$
    \item $\displaystyle \int\frac{1}{\tan(x/2)+\cot(x/2)}dx$
  \end{enumerate}

\item 다음 물음에 답하여라.
\begin{enumerate}[label=(\arabic*)]
    \item $x=a\sin\theta(-\dfrac{\pi}{2}\le\theta\le\dfrac{\pi}{2})$ 로 치환하여 $\int\dfrac{1}{\sqrt{(a^{2}-x^{2})^{3}}}dx(a>0)$ 를 구하여라.
    \item $x=a\tan\theta(-\dfrac{\pi}{2}<\theta<\dfrac{\pi}{2})$ 로 치환하여 $\int\dfrac{1}{\sqrt{(a^{2}+x^{2})^{3}}}dx(a>0)$ 를 구하여라.
\end{enumerate}

\item $f(x)=\int(x \ln x+e^{x}+x+1)dx$ 일 때, 다음 극한값을 구하여라.
    $\displaystyle \lim_{h\rightarrow0}\frac{f(1+h+h^{2})-f(1-h)}{h}$

\item 미분가능한 함수 $f(x)$가 모든 실수 $x$에 대하여 
\[xf'(x)-f(x)=(x^4+2x^2)e^(x-1)\]을 만족시킨다. $f(1)=3$일때, $f(3)$의 값을 구하여라.

\item 함수 $f(x), g(x)$ 가 다음 세 조건을 만족시킬 때, $f(x)+g(x)$ 를 구하여라.
\begin{enumerate}[label=(\text{\roman*}), leftmargin=10mm, nosep]
    \item $f^{\prime}(x)=2g(x), \ g^{\prime}(x)=2f(x)$
    \item $f(0)=1, \ g(0)=e-1$
    \item $f(x)>0, \ g(x)>0$    
\end{enumerate}

\end{enumerate}


\end{document}